\section{构造低次替代型}\label{sec:construction}
计数例程对每对 $(f,v)$（因子 $f$ 读取变量 $v$）需要一个替代型。
对第~\ref{sec:setup} 节的应用，$f$ 是一个多项式方程，
其 scope 是其各单项变量集的并。本节只使用这些 scope。
我们构造定义~\ref{def:layout} 中的行 $\vec{P}_{f,v}(T)$
与测量矩阵 $C_{t,g}$，连同划分 $S_{t,g}$ 与因子指派
$\sigma_f$。主要困难在于保持 $\deg \vec{P}_{f,v}\le R$
且 $R$ 远小于 $Q$。$H$ 的次数界随 $R$ 线性增长，
因此小的 $R/Q$ 减少每点所需的 Taylor 系数个数。
在 $T=0,1,\ldots,Q$ 处独立选行再插值，一般给出次数 $Q$。
我们改为把它们在所有这些点上的取值一起构造出来。

构造先在辅助坐标的完整阵列上给出多项式行。
随后把每个「因子--变量」对放入该阵列的一个副本中。
多个位置可以含同一布尔变量 $x_v$；其取值总相等，
在计数和中不会成为独立变量。

\subsection{辅助坐标}\label{subsec:incidence-layout}
我们先指定更大的坐标集 $\mathcal B$——辅助行
$\vec{P}_u(T)$ 将作用其上——以及把实际的 $(f,v)$
放入其副本所需的关联界。
设 $(I_f)_{f\in\mathcal F}$ 为因子 scope，$I_f\subseteq[n]$，
写
\begin{equation}\label{eq:factor-incidences}
 \begin{gathered}
 \mathcal A=\{(f,v):f\in\mathcal F,\ v\in I_f\},\qquad M=|\mathcal A|,\\
 |I_f|\le w\quad(f\in\mathcal F),\qquad
 \#\{f:v\in I_f\}\le b\quad(v\in[n]).
 \end{gathered}
\end{equation}
于是 $w$ 界住一个因子 scope 中的变量数，$b$
界住含一个变量的因子数。scope 界 $w$ 不必界定因子多项式
的次数。scope 为空的因子是常数，不需要辅助位置；
在每个点把这样的因子指派到组一。
如有必要，对每个 $v\in[n]$ 附加一个 scope 为 $\{v\}$
的互异常数一因子。这既不改变乘积也不改变计数，
且使最大关联至多增加一。以 $b$ 记此次添加后的界，我们设
\[
 n\ge1,\qquad b,w\ge1,\qquad n\le M\le bn.
\]
本节始终设 $r\ge2$、$1\le R<Q$ 为整数，$0<\tau\le1/2$。

在每个点 $t\in[Q]$，变量 $x_v$ 将属于一个组，
而含 $v$ 的因子可能被指派到另一个组。我们将选取
\[
 a_v\in[r]^Q\quad(v\in[n]),\qquad
 d_f\in[r]^Q\quad(f\in\mathcal F),
\]
其中 $a_v(t)$ 指定含 $x_v$ 的组，$d_f(t)$ 指定指派给因子
$f$ 的组。位置 $(f,v)$ 将有标号 $(a_v,d_f)$。
特别地，所有含 $x_v$ 的位置都有同一个第一序列 $a_v$。

在选取 $a_v$ 与 $d_f$ 之前，先取全部可能标号的完整集合：
\begin{equation}\label{eq:auxiliary-labels}
 \mathcal B=([r]^Q)^2,\qquad u=(a_u,d_u),\qquad D=|\mathcal B|=r^{2Q}.
\end{equation}
辅助向量 $\zeta\in\mathbb Q^{\mathcal B}$
为每个对 $u$ 有一个有理坐标 $\zeta_u$。
标号 $u$ 在 $\zeta$ 变动时保持固定。
对 $t\in[Q]$ 与 $g\in[r]$，定义
\begin{equation}\label{eq:auxiliary-groups}
 S^{\mathcal B}_{t,g}=\{u\in\mathcal B:a_u(t)=g\}.
\end{equation}
对固定的 $t$，这些集合把 $\mathcal B$ 划分为 $r$ 个大小
$D/r$ 的集合。

对每个 $u\in\mathcal B$，我们寻求次数至多 $R-1$、
满足 $\vec{P}_u(0)=\delta_u$ 的多项式行
$\vec{P}_u(T)\in\mathbb Q[T]^{1\times\mathcal B}$。
在 $T=t$，取值 $\langle \vec{P}_u(t), \zeta\rangle$
应由
\[
 \zeta|_{S^{\mathcal B}_{t,d_u(t)}}
 \quad\text{与}\quad C^{\mathcal B}\zeta
\]
确定，其中 $C^{\mathcal B}$ 是一个与 $u,t$ 无关的整数矩阵，
分量在 $\{0,1,-1\}$ 中且行数很少。
稍后使用 $u_{f,v}=(a_v,d_f)$ 将使因子 $f$
的每个输入都用组 $d_f(t)$。$a_v,d_f$ 的选取与放入副本
分别推迟到第~\ref{subsec:choosing-labels}
与~\ref{subsec:embedding} 节。

\subsection{矩阵 $M_0$ 与 $C^{\mathcal B}$}
为构造行 $\vec{P}_u(T)$，我们寻求固定 $\ker C^{\mathcal B}$
并满足
\begin{equation}\label{eq:matrix-support-target}
 M_0[u,y]\ne0\quad\Longrightarrow\quad
 \bigl|\{s\in[Q]:a_y(s)\ne d_u(s)\}\bigr|\le R-1.
\end{equation}
的矩阵 $M_0\in\mathbb Q^{\mathcal B\times\mathcal B}$。
恒等式
\begin{equation}\label{eq:coordinate-decomposition}
 \zeta_u=((I-M_0)\zeta)_u+
          \sum_{y\in\mathcal B}M_0[u,y]\zeta_y
\end{equation}
解释了这两个要求。若 $M_0$ 固定 $\ker C^{\mathcal B}$，
则第一项由 $C^{\mathcal B}\zeta$ 确定；我们将使它与 $T$
无关。在和式中，把每个非零 $M_0[u,y]$ 乘以一个在零点为一、
在每个满足 $a_y(t)\ne d_u(t)$ 的 $t$ 处为零的多项式。
界 \eqref{eq:matrix-support-target} 容许次数至多 $R-1$。
于是在 $T=t$，和式将只用坐标 $y\in S^{\mathcal B}_{t,d_u(t)}$。

对 $s\in[Q]$，设 $\vartheta_s:\mathcal B\to\mathcal B$
交换 $a_u(s)$ 与 $d_u(s)$，其余分量不变。
用
\[
 (\Pi_s\zeta)_u=\zeta_{\vartheta_s(u)}.
\]
定义其置换矩阵 $\Pi_s$。
被交换的对 $\vartheta_s(u)$ 属于 $\mathcal B$，
因为该集合含 $[r]^Q$ 中序列的每一对。
矩阵 $\Pi_s$ 两两交换且满足 $\Pi_s^2=I$。其乘积
\[
 V_0=\prod_{s=1}^Q\Pi_s
\]
交换完整序列 $a_u$ 与 $d_u$。于是
$V_0[u,y]\ne0$ 蕴含 $a_y(s)=d_u(s)$ 对每个 $s$；
它以零错配满足 \eqref{eq:matrix-support-target}。
我们将修改 $V_0$，使所得矩阵同时固定一列短的整数测量的核。

定义
\begin{equation}\label{eq:swap-weight}
 W_0=\frac12\sum_{s=1}^Q(I-\Pi_s).
\end{equation}
在 $\Pi_s$ 的公共特征向量 $\psi$ 上，设 $j$ 为满足
$\Pi_s\psi=-\psi$ 的下标 $s$ 的个数。则
\[
 W_0\psi=j\psi,\qquad V_0\psi=(-1)^j\psi,
 \qquad V_0p(W_0)\psi=(-1)^jp(j)\psi
\]
对每个多项式 $p$ 成立。因此只要 $p(j)=(-1)^j$，
$V_0p(W_0)$ 就固定 $\psi$。
我们将构造一个显式公共特征基，其中具有给定 $j$
的向量比例恰为 $\Pr[W=j]$
（$W\sim\operatorname{Bin}(Q,(r-1)/(2r))$）。

对区间
$\mathcal I=\{A_0,\ldots,A_0+R-1\}\subseteq\{0,\ldots,Q\}$，
可以用次数至多 $R-1$ 的多项式强制等式
$f_{\mathcal I}(j)=(-1)^j$ 对 $j\in\mathcal I$ 成立。
我们使用显式插值公式
\begin{equation}\label{eq:swap-matrix-definition}
 f_{\mathcal I}(z)=(-1)^{A_0}\sum_{j=0}^{R-1}(-2)^j
                         \binom{z-A_0}{j},
 \qquad M_0=V_0f_{\mathcal I}(W_0),
\end{equation}
其中
\[
 \binom{z-A_0}{j}=\frac1{j!}\prod_{a=0}^{j-1}(z-A_0-a)
\]
且空乘积为一。对 $0\le q<R$，二项式定理给出
$f_{\mathcal I}(A_0+q)=(-1)^{A_0}(1-2)^q=(-1)^{A_0+q}$。

\begin{lemma}\label{lem:swaps-matrix}
对如上的 $\mathcal B$、$\mathcal I$ 与 $M_0$，
存在矩阵 $C^{\mathcal B}\in\{0,1,-1\}^{[c_{\mathcal B}]\times\mathcal B}$，
使
\[
 c_{\mathcal B}\le D\Pr[W\notin\mathcal I],\qquad
 W\sim\operatorname{Bin}\!\left(Q,\frac{r-1}{2r}\right),
\]
且对每个 $\xi\in\ker C^{\mathcal B}$ 有
$M_0\xi=\xi$。
\end{lemma}
\begin{proof}
对 $(\alpha,\delta)\in[r]^2$，设 $e_{\alpha\delta}$ 为该
有序对的指示函数。交换其两个分量有正交特征基
\[
 e_{\alpha\alpha}\quad(\alpha\in[r]),\qquad
 e_{\alpha\delta}+e_{\delta\alpha},\quad
 e_{\alpha\delta}-e_{\delta\alpha}\quad(1\le\alpha<\delta\le r).
\]
前两类特征值为 $+1$，第三类为 $-1$。
特征值 $+1$ 的向量有 $r(r+1)/2$ 个，$-1$ 的有
$r(r-1)/2$ 个。

对每个 $s\in[Q]$，选取这些函数之一
$\psi_s:[r]^2\to\{0,1,-1\}$，并定义
\[
 \psi(u)=\prod_{s=1}^Q\psi_s(a_u(s),d_u(s))
 \qquad(u\in\mathcal B).
\]
$D=r^{2Q}$ 种选取给出非零正交向量：其内积分解为各
$\psi_s$ 内积的乘积。它们构成 $\mathbb Q^{\mathcal B}$
的一组分量在 $\{0,1,-1\}$ 中的基。
若恰有 $j$ 个所选 $\psi_s$ 的特征值为 $-1$，则
$W_0\psi=j\psi$、$V_0\psi=(-1)^j\psi$。
故只要 $j\in\mathcal I$，$M_0$ 就固定 $\psi$。

为使 $\ker C^{\mathcal B}$ 成为被 $M_0$ 固定的空间，
取 $C^{\mathcal B}$ 的行为 $j\notin\mathcal I$
的那些基向量的转置。每个位置上有比例 $(r-1)/(2r)$
的选取特征值为 $-1$，且各位置独立。
因此作为行保留的 $D$ 个基向量的比例是
$\Pr[W\notin\mathcal I]$，证得行数界。
若 $C^{\mathcal B}\xi=\vec{0}$，则 $\xi$
与每个保留的基向量正交，故是 $j\in\mathcal I$
的那些基向量的线性组合，而它们全部被 $M_0$ 固定。
\end{proof}

对界 $\epsilon>0$，条件
\begin{equation}\label{eq:incidence-tail}
 \Pr\!\left[\operatorname{Bin}\!\left(Q,\frac{r-1}{2r}\right)
                         \notin\mathcal I\right]\le\epsilon
\end{equation}
给出 $c_{\mathcal B}\le\epsilon D$。
放置因子输入后，我们将取 $\epsilon=\tau/(rb)$。
接下来证明 $M_0$ 的另一个要求。

\begin{lemma}\label{lem:swaps-support}
对所有 $u,y\in\mathcal B$，若 $M_0[u,y]\ne0$，则
\[
 \bigl|\{s\in[Q]:a_y(s)\ne d_u(s)\}\bigr|\le R-1.
\]
\end{lemma}
\begin{proof}
在两两交换的矩阵 $\Pi_s$ 中展开 $f_{\mathcal I}(W_0)$。
由 $\deg f_{\mathcal I}\le R-1$ 与 $\Pi_s^2=I$，
每个所得单项是某个满足 $|J|\le R-1$ 的
$J\subseteq[Q]$ 对应的 $\prod_{s\in J}\Pi_s$
的标量倍。乘以 $V_0$ 给出
\[
 V_0\prod_{s\in J}\Pi_s=\prod_{s\notin J}\Pi_s.
\]
该置换矩阵的 $(u,y)$ 分量非零要求
$a_y(s)=d_u(s)$ 对每个 $s\notin J$。
其余位置至多 $R-1$ 个。若不一致的位置多于 $R-1$ 个，
则每个被加项的 $(u,y)$ 分量为零，
故 $M_0[u,y]=0$。
\end{proof}

\subsection{多项式行}
现在用 $M_0$ 构造第~\ref{subsec:incidence-layout} 节
所求的辅助行 $\vec{P}_u(T)$。
引理~\ref{lem:swaps-matrix} 使 $(I-M_0)\zeta$
成为 $C^{\mathcal B}\zeta$ 的函数。
引理~\ref{lem:swaps-support} 界住非零项
$M_0[u,y]\zeta_y$ 使用 $S^{\mathcal B}_{t,d_u(t)}$
之外坐标的点数。现在把该项乘以在这些点为零、在零点为
一的多项式。

对 $M_0[u,y]\ne0$，定义
\begin{equation}\label{eq:coefficient-correction}
 E_{u,y}=\{s\in[Q]:a_y(s)\ne d_u(s)\},\qquad
 \varphi_{u,y}(T)=\prod_{s\in E_{u,y}}\left(1-\frac Ts\right).
\end{equation}
空乘积为一。则 $\deg\varphi_{u,y}\le R-1$、
$\varphi_{u,y}(0)=1$，且只要 $a_y(t)\ne d_u(t)$ 就有
$\varphi_{u,y}(t)=0$。定义
\begin{equation}\label{eq:swaps-row}
 \vec{P}_u(T)=\delta_u(I-M_0)+
 \sum_{\substack{y\in\mathcal B\\M_0[u,y]\ne0}}
 M_0[u,y]\varphi_{u,y}(T)\delta_y.
\end{equation}
该行作为 $\zeta$ 的线性型作用，其系数是单个参数 $T$
的多项式。

\begin{lemma}\label{lem:swaps-forms}
\eqref{eq:swaps-row} 中的行次数至多 $R-1$，
且对每个 $u\in\mathcal B$ 满足 $\vec{P}_u(0)=\delta_u$。
对 $t\in[Q]$ 与 $g\in[r]$，设 $C^{\mathcal B}_{t,g}$
为 $C^{\mathcal B}$ 限制到列 $S^{\mathcal B}_{t,g}$。
对每个满足 $g\ne d_u(t)$ 的 $u,t,g$ 与每个
$\xi\in\mathbb Q^{S^{\mathcal B}_{t,g}}$，其补零延拓
$\widetilde\xi$ 满足
\[
 C^{\mathcal B}_{t,g}\xi=\vec{0}
 \quad\Longrightarrow\quad \langle \vec{P}_u(t), \widetilde\xi\rangle=0.
\]
更一般地，$\langle \vec{P}_u(t), \zeta\rangle$
由 $\zeta|_{S^{\mathcal B}_{t,d_u(t)}}$ 与
$C^{\mathcal B}\zeta$ 确定。
\end{lemma}
\begin{proof}
次数界来自 $\varphi_{u,y}$ 的次数界。在零点，
\[
 \vec{P}_u(0)=\delta_u(I-M_0)+\delta_uM_0=\delta_u.
\]
把 \eqref{eq:swaps-row} 作用于 $\zeta$ 给出
\begin{equation}\label{eq:swaps-form-decomposition}
 \langle \vec{P}_u(T), \zeta\rangle=((I-M_0)\zeta)_u+
 \sum_{\substack{y\in\mathcal B\\M_0[u,y]\ne0}}
 M_0[u,y]\varphi_{u,y}(T)\zeta_y.
\end{equation}
在 $T=t$，所有满足 $a_y(t)\ne d_u(t)$ 的被加项为零。
若 $\widetilde\xi$ 支撑在 $S^{\mathcal B}_{t,g}$
（$g\ne d_u(t)$）上，则每个剩余被加项在 $\widetilde\xi$
上为零。假设 $C^{\mathcal B}_{t,g}\xi=\vec{0}$
也给出 $C^{\mathcal B}\widetilde\xi=\vec{0}$，
故由引理~\ref{lem:swaps-matrix} 有
$(I-M_0)\widetilde\xi=\vec{0}$。
这就证明了该蕴含。

对最后的断言，在 $C^{\mathcal B}$ 下测量相同的两个向量
相差其核的一个元素。于是 $(I-M_0)\zeta$
由 $C^{\mathcal B}\zeta$ 确定。
\eqref{eq:swaps-form-decomposition} 中在 $T=t$
的其余项只用 $\zeta|_{S^{\mathcal B}_{t,d_u(t)}}$。
\end{proof}

固定 $t$ 并记 $\zeta_g=\zeta|_{S^{\mathcal B}_{t,g}}$。
$C^{\mathcal B}$ 的限制满足
\begin{equation}\label{eq:auxiliary-measurements}
 C^{\mathcal B}\zeta=\sum_{g=1}^r C^{\mathcal B}_{t,g}\zeta_g.
\end{equation}
因此固定每个向量 $C^{\mathcal B}_{t,g}\zeta_g$
既固定 $C^{\mathcal B}\zeta$，又对 $r$ 个辅助坐标组
施加各自的条件。当全部 $C^{\mathcal B}_{t,g}\zeta_g$
固定时，$\langle \vec{P}_u(t), \zeta\rangle$
由 $\zeta|_{S^{\mathcal B}_{t,d_u(t)}}$ 确定。
把 $(f,v)$ 放到标号 $(a_v,d_f)$ 的坐标处，
我们将在实际布尔变量上得到相应的分离。

\begin{remark}
我们演示 \eqref{eq:swaps-row} 如何把指派组之外的依赖
限制到一个线性测量。取 $Q=R=2$（在运行时间界所用
$R<Q$ 体制之外）与 $r=2$。
用 $\{0,1\}$ 作组下标，取 $\mathcal I=\{0,1\}$，
故 $f_{\mathcal I}(z)=1-2z$。
从标号 $(00,11)$ 出发，两次交换生成四个坐标：
\[
 \zeta_0=\zeta_{(00,11)},\quad \zeta_1=\zeta_{(10,01)},\quad
 \zeta_2=\zeta_{(01,10)},\quad \zeta_{12}=\zeta_{(11,00)}.
\]
两次交换都保持这四个坐标，故可把计算限制在其上。
在此限制上，$M_0=\Pi_1+\Pi_2-\Pi_1\Pi_2$
固定
\[
 c(\zeta)=\zeta_0-\zeta_1-\zeta_2+\zeta_{12}
\]
的核。以 $u=(00,11)$ 为下标的行给出
\[
 \langle \vec{P}_u(T), \zeta\rangle=\zeta_0-\frac T2\zeta_1-T\zeta_2.
\]
它在零点等于 $\zeta_0$，而
\[
 \langle \vec{P}_u(1), \zeta\rangle=c(\zeta)+\frac12\zeta_1-\zeta_{12},\qquad
 \langle \vec{P}_u(2), \zeta\rangle=c(\zeta)-\zeta_2-\zeta_{12}.
\]
在 $t=1$，$\zeta_1,\zeta_{12}$ 的标号满足
$a_y(1)=1=d_u(1)$。在 $t=2$，$\zeta_2,\zeta_{12}$
的标号满足 $a_y(2)=1=d_u(2)$。
因此固定 $c(\zeta)$ 后，$\langle \vec{P}_u(t), \zeta\rangle$
只依赖其指派组中的坐标。
这里 $c(\zeta)$ 在四坐标限制上扮演 $C^{\mathcal B}\zeta$
的角色。完整构造必须对每个 $u\in\mathcal B$
以至多 $\epsilon D$ 个测量行给出该性质。
\end{remark}

\subsection{选取变量组与因子组}\label{subsec:choosing-labels}
现在从第~\ref{subsec:incidence-layout} 节选取族
$(a_v)_{v\in[n]}$ 与 $(d_f)_{f\in\mathcal F}$，并置
\[
 u_{f,v}=(a_v,d_f)\qquad((f,v)\in\mathcal A).
\]
这些选取有两个要求。第一，若某标号出现 $c_u$ 次，
其关联在辅助阵列副本中需要 $c_u$ 个不同位置。
因此所需副本数为
\[
 c_u=\#\{(f,v)\in\mathcal A:u_{f,v}=u\},\qquad
 h=\max_{u\in\mathcal B}c_u.
\]
第二，实际变量组 $\{v:a_v(t)=g\}$ 必须小。
只界定 $h$ 界住的是一组内的关联个数，
而非其布尔变量个数。下面的引理同时控制两个量。

\begin{lemma}\label{lem:incidence-labels}
设 $(I_f)_{f\in\mathcal F}$ 满足
\eqref{eq:factor-incidences}，$n\le M\le bn$，
并对整数 $r\ge2$、$Q\ge1$ 置 $D=r^{2Q}$。设
\begin{equation}\label{eq:incidence-balance-hypotheses}
 2(b+w)D^2\le M,\qquad 2Qr^2\le n.
\end{equation}
可以确定性选取 $a_v,d_f\in[r]^Q$
（对 scope 为空的因子取 $d_f(t)=1$），使
\[
 \begin{aligned}
 \max_{u\in\mathcal B}\#\{(f,v)\in\mathcal A:(a_v,d_f)=u\}
 &\le\frac{2M}{D},\\
 \#\{v\in[n]:a_v(t)=g\}&\le\frac{2n}{r}
 \qquad(t\in[Q],\ g\in[r]).
 \end{aligned}
\]
构造花费 $\poly(n,M,|\mathcal F|,D,Q,r,b,w)$ 次位运算。
\end{lemma}
\begin{proof}
我们先估计随机标号的效果，然后再确定性固定其取值。
在此估计中，所有 $a_v$ 与所有 $I_f$ 非空的 $d_f$
独立均匀取自 $[r]^Q$。对空的 $I_f$，恒置 $d_f$ 为一。
每个关联标号 $u_{f,v}=(a_v,d_f)$ 均匀分布于 $\mathcal B$。
置 $n_{t,g}=\#\{v:a_v(t)=g\}$，考虑
\[
 \Phi=\sum_{u\in\mathcal B}(c_u-M/D)^2,
 \qquad
 \Xi=\sum_{t=1}^Q\sum_{g=1}^r(n_{t,g}-n/r)^2.
\]
其作为指示子和的展开为
\[
 \begin{aligned}
 \Phi&=\sum_{(f,v),(f',v')\in\mathcal A}
       \mathds{1}_{\{(a_v,d_f)=(a_{v'},d_{f'})\}}-\frac{M^2}{D},\\
 \Xi&=\sum_{t=1}^Q\sum_{v,v'\in[n]}
       \mathds{1}_{\{a_v(t)=a_{v'}(t)\}}-\frac{Qn^2}{r}.
 \end{aligned}
\]
两个关联标号除非共享因子或变量，否则独立。
对每个 $(f,v)$，至多有 $b+w$ 个关联共享 $f$ 或 $v$
（含 $(f,v)$ 自身）。独立对的相等概率为 $1/D$；
其余对的概率至多为一。故
\[
 \mathbb E\Phi\le M(b+w).
\]
对每个 $t$ 与 $v\ne v'$，$a_v(t),a_{v'}(t)$
独立均匀取自 $[r]$。第二个展开因此给出
\[
 \mathbb E\Xi=Qn(1-1/r)\le Qn.
\]

为确定性获得标号，我们逐个固定 $a_v$ 与非空 $I_f$
的 $d_f$。每步选取一个不增加
\[
 \frac{\Phi}{M(b+w)}+\frac{\Xi}{Qn}
\]
的条件期望的取值。这样的选取存在，因为当前期望是下一个
$a_v$ 或 $d_f$ 各可能取值上的平均。
最终值至多为一的两倍。由于两个被加项都非负，
\[
 c_u\le\frac MD+\sqrt{2M(b+w)},\qquad
 n_{t,g}\le\frac nr+\sqrt{2Qn}.
\]
\eqref{eq:incidence-balance-hypotheses} 的两个不等式
使两个根号项分别至多 $M/D$ 与 $n/r$。

下面的枚举在 $\poly(n,M,|\mathcal F|,D,Q,r,b,w)$
次位运算内计算条件期望。
置 $K=r^Q$，则 $D=K^2$。$\Phi$ 展开中的每个指示子至多
依赖四个独立选取 $a_v,d_f,a_{v'},d_{f'}$；
重复的下标只贡献一个选取。$\Xi$ 中的每个指示子至多依赖
$a_v,a_{v'}$。在任何部分固定之后，对一个指示子中剩余
选取枚举至多 $K^4=D^2$ 个赋值。其条件概率的分母整除
$K^4$。对 $M^2+Qn^2$ 个指示子求和即算出两个期望。
计入归一化 $M(b+w)$ 与 $Qn$，这些计算中的全部有理值的
分子与分母都有
\[
 O\!\left(\log(2D)+\log(2M)+\log(2n)+\log(2Q)+\log(2+b+w)\right)
\]
位。需要固定的选取至多 $n+|\mathcal F|$ 个，
每个有 $K$ 个候选值。由 $K^4=D^2$，
整个构造花费 $\poly(n,M,|\mathcal F|,D,Q,r,b,w)$
次位运算。
\end{proof}

\subsection{把变量副本放入辅助阵列}\label{subsec:embedding}
现在把辅助行 $\vec{P}_u(T)$ 与矩阵 $C^{\mathcal B}$
转移到定义~\ref{def:layout} 所要求的、原变量上的行
$\vec{P}_{f,v}(T)$ 与矩阵 $C_{t,g}$。
设 \eqref{eq:incidence-balance-hypotheses} 成立并应用
引理~\ref{lem:incidence-labels}。
以 $h=\max_u c_u$，用 $[h]\times\mathcal B$
为辅助坐标的 $h$ 份副本编下标。
对每个 $(f,v)\in\mathcal A$，指派 $c(f,v)\in[h]$
使所有对 $(c(f,v),u_{f,v})$ 互异。
这是可能的，因为没有标号出现超过 $h$ 次，且
\[
 hD\le2M\le2bn.
\]
定义线性映射
$\iota:\mathbb Q^n\longrightarrow\mathbb Q^{[h]\times\mathcal B}$：
\[
 (\iota x)_{c,u}=
 \begin{cases}
 x_v,&(c,u)=(c(f,v),u_{f,v})\text{ 对某个 }(f,v)\in\mathcal A,\\
 0,&\text{否则}.
 \end{cases}
\]
被占据的位置互异，故该定义无歧义。
不同位置可以含同一个 $x_v$，未用的位置固定为零。
对 $c\in[h]$，定义投影
\[
 \pi_c:\mathbb Q^{[h]\times\mathcal B}\longrightarrow\mathbb Q^{\mathcal B},
 \qquad (\pi_c\eta)_u=\eta_{c,u}.
\]
现在置
\begin{equation}\label{eq:actual-rows}
 \begin{gathered}
 \vec{P}_{f,v}(T)=\vec{P}_{u_{f,v}}(T)\pi_{c(f,v)}\iota,\\
 S_{t,g}=\{v\in[n]:a_v(t)=g\},\qquad \sigma_f(t)=d_f(t).
 \end{gathered}
\end{equation}
行 $\vec{P}_{f,v}(T)$ 作用在实际变量 $x\in\mathbb Q^n$
上。在 $T=0$，它提取位置 $(c(f,v),u_{f,v})$
处的值，即 $x_v$。于是 $\vec{P}_{f,v}(0)=\delta_v$，
次数至多 $R-1$。引理~\ref{lem:incidence-labels}
给出 $|S_{t,g}|\le2n/r$。

用
\begin{equation}\label{eq:actual-measurements}
 \widetilde C=\operatorname{diag}
   \bigl(\underbrace{C^{\mathcal B},\ldots,C^{\mathcal B}}_{h\text{ 份副本}}\bigr),
 \qquad \widehat C=\widetilde C\iota,
 \qquad C_{t,g}=\widehat C|_{S_{t,g}}.
\end{equation}
定义测量矩阵。
该限制保留 $S_{t,g}$ 中的列与全部行（含零行）。
于是 $\widehat Cx$ 在每个副本的坐标从 $x$ 填充后
记录其辅助测量。
矩阵 $C_{t,g}$ 是整数的。其分量不必属于 $\{0,1,-1\}$，
因为含同一 $x_v$ 的多个位置上的系数被相加。
$\widetilde C$ 的每行至多 $D$ 个非零分量，
全为一或 $-1$。乘以 $\iota$ 在公共列中添加分量并丢弃
未用位置，故不会增加其绝对值之和。
因此 $C_{t,g}$ 的每行绝对行和至多 $D$，
其在布尔赋值上的值落在 $[-D,D]\cap\mathbb Z$ 中。

若 \eqref{eq:incidence-tail} 成立，则每个 $C_{t,g}$
有 $c_{t,g}=hc_{\mathcal B}$ 行。故
\begin{equation}\label{eq:incidence-measurement-count}
 \sum_{g=1}^r c_{t,g}=rh c_{\mathcal B}
 \le rh\epsilon D\le2r\epsilon M\le2r\epsilon bn.
\end{equation}
取 $\epsilon=\tau/(rb)$ 使其至多 $2\tau n$。

为验证分离性，固定 $t\in[Q]$、$g\in[r]$ 与满足
$g\ne\sigma_f(t)$ 的 $(f,v)\in\mathcal A$。
设 $\xi\in\mathbb Q^{S_{t,g}}$ 满足
$C_{t,g}\xi=\vec{0}$，补零延拓为
$\widetilde\xi\in\mathbb Q^n$。
\eqref{eq:actual-measurements} 给出
$\widetilde C\iota\widetilde\xi=\vec{0}$。
于是对每个 $c\in[h]$，
\[
 C^{\mathcal B}\pi_c\iota\widetilde\xi=\vec{0}.
\]
$\pi_c\iota\widetilde\xi$ 的每个非零坐标含一个满足
$a_{v'}(t)=g$ 的变量 $v'$。其辅助标号的第一序列为
$a_{v'}$，故该向量支撑在 $S^{\mathcal B}_{t,g}$ 上。
在副本 $c(f,v)$ 中以 $d_{u_{f,v}}(t)=d_f(t)=\sigma_f(t)$
应用引理~\ref{lem:swaps-forms}，给出
\[
 \langle \vec{P}_{f,v}(t), \widetilde\xi\rangle
 =\langle \vec{P}_{u_{f,v}}(t), \pi_{c(f,v)}\iota\widetilde\xi\rangle=0.
\]
这正是定义~\ref{def:layout} 中的核条件。
因子 $f$ 的全部输入使用同一指派组 $d_f(t)$，
而计数和仍恰遍历 $n$ 个布尔位。

\subsection{系数界与完成的构造}
\eqref{eq:actual-rows} 中的代入行
$\vec{P}_{f,v}(T)\in\mathbb Q[T]^{1\times n}$
的分量是带理系数的多项式。现在界定存储这些系数所需的
位数；这些界将在第~\ref{subsec:one-query-cost} 节
控制精确算术。我们先对 \eqref{eq:swaps-row} 定义的
辅助行 $\vec{P}_u(T)$ 给出公分母与系数界。
由 \eqref{eq:actual-rows} 从 $\vec{P}_u(T)$
转到 $\vec{P}_{f,v}(T)$ 只在公共列中添加系数或丢弃
未用列，故同一分母与绝对值之和的界仍然适用。

\begin{lemma}\label{lem:row-coefficients}
设 $d_{\mathrm{den}}=(R-1)!Q!$。则对每个 $u\in\mathcal B$
有 $d_{\mathrm{den}}\vec{P}_u(T)\in\mathbb Z[T]^{1\times D}$。
行 $d_{\mathrm{den}}\vec{P}_u(T)$ 中全部系数的绝对值之
和至多 $2^{O(Q\log Q)}$。对 $d_{\mathrm{den}}\vec{P}_{f,v}(T)$
同样成立。矩阵 $M_0,C^{\mathcal B}$ 与行 $\vec{P}_u(T)$
可在 $\poly(D,Q,r)$ 次位运算内构造。
\end{lemma}
\begin{proof}
对 $0\le j<R$，$f_{\mathcal I}(W_0)$ 的相应被加项为
\[
 (-1)^{A_0}(-2)^j\binom{W_0-A_0I}{j}
 =\frac{(-1)^{A_0+j}}{j!}
   \prod_{a=0}^{j-1}\bigl(2W_0-2(A_0+a)I\bigr).
\]
由 $2W_0=QI-\sum_s\Pi_s$，右端每个矩阵因子都是整数的。
且 $j!$ 整除 $(R-1)!$。
因此乘以置换矩阵 $V_0$ 表明 $(R-1)!M_0$ 是整数的。
$\varphi_{u,y}$ 的系数的分母整除
$\prod_{s\in E_{u,y}}s$，而后者整除 $Q!$。
\eqref{eq:swaps-row} 证明所述公分母。

对系数估计，用矩阵范数
$\|A\|_\infty=\max_i\sum_j|A[i,j]|$
与多项式范数 $\|\sum_j a_jT^j\|_1=\sum_j|a_j|$。
由 $\|W_0\|_\infty\le Q$ 与上式各因子满足
$0\le A_0+a\le Q$，每个 $2W_0-2(A_0+a)I$
的范数至多 $4Q$。于是
\[
 \|M_0\|_\infty\le R(4Q)^{R-1},\qquad
 \|\varphi_{u,y}\|_1
 \le\prod_{s\in E_{u,y}}(1+1/s)\le2^Q.
\]
行 $d_{\mathrm{den}}\vec{P}_u$ 中全部系数的绝对值之
和因此至多
\[
 d_{\mathrm{den}}\bigl(1+(1+2^Q)R(4Q)^{R-1}\bigr).
\]
由 $R<Q$，其对数为 $O(Q\log Q)$。
对 $u=u_{f,v}$，把 $d_{\mathrm{den}}\vec{P}_u(T)$
乘以 $\pi_{c(f,v)}\iota$ 在公共列中添加系数并丢弃未用列。
由 \eqref{eq:actual-rows}，所得行是
$d_{\mathrm{den}}\vec{P}_{f,v}(T)$，故满足同一界。

枚举引理~\ref{lem:swaps-matrix} 的特征基、求值
$f_{\mathcal I}(W_0)$、并作成多项式 $\varphi_{u,y}$
需要 $\poly(D,Q,r)$ 次算术运算。使用上述公分母，
这些矩阵与多项式计算中存储的系数的分子与分母都有
$O(Q\log Q)$ 位。因此构造花费 $\poly(D,Q,r)$
次位运算。
\end{proof}

构造有三个定量要求。
\eqref{eq:incidence-balance-hypotheses} 的两个不等式
界定辅助副本数与每组实际变量数。
取 $\epsilon=\tau/(rb)$ 的尾界经由
\eqref{eq:incidence-measurement-count} 给出总测量界。
我们汇总结果供计数算法使用。

\begin{lemma}\label{lem:layout}
设整数 $n,b,w\ge1$，$(I_f)_{f\in\mathcal F}$
为 $[n]$ 的规模至多 $w$ 的子集，
每个 $v\in[n]$ 至多属于 $b$ 个集合。
置 $M=\sum_f|I_f|$ 并设 $n\le M\le bn$。
设 $r\ge2$、$1\le R<Q$ 为整数，$0<\tau\le1/2$，
$D=r^{2Q}$。设
\[
 2(b+w)D^2\le M,\qquad 2Qr^2\le n,
\]
且区间
$\mathcal I=\{A_0,\ldots,A_0+R-1\}\subseteq\{0,\ldots,Q\}$
满足
\[
 \Pr[W\notin\mathcal I]\le\frac{\tau}{rb},\qquad
 W\sim\operatorname{Bin}\!\left(Q,\frac{r-1}{2r}\right).
\]
则可以在 $\poly(n,M,|\mathcal F|,D,Q,r,b,w)$ 次位运算内
确定性构造定义~\ref{def:layout} 中的全部数据。
\end{lemma}
\begin{proof}
引理~\ref{lem:swaps-matrix} 构造满足
$c_{\mathcal B}\le\tau D/(rb)$ 的 $C^{\mathcal B}$。
引理~\ref{lem:swaps-support} 与~\ref{lem:swaps-forms}
给出辅助多项式行。
引理~\ref{lem:incidence-labels} 选取变量组与因子组，
满足 $hD\le2M$ 与 $|S_{t,g}|\le2n/r$。
\eqref{eq:actual-rows} 与 \eqref{eq:actual-measurements}
中的映射随后给出满足 $\vec{P}_{f,v}(0)=\delta_v$、
次数至多 $R-1\le R$ 与所需核蕴含的行。
$C_{t,g}$ 的每行绝对行和至多 $D$，且
\eqref{eq:incidence-measurement-count} 对每个 $t\in[Q]$
给出 $\sum_{g=1}^r c_{t,g}\le2\tau n$。
引理~\ref{lem:row-coefficients} 提供公分母与系数界。
合起来，各构造花费 $\poly(n,M,|\mathcal F|,D,Q,r,b,w)$
次位运算，所得阵列占用
$\poly(n,M,|\mathcal F|,D,Q,r,b,w)$ 比特。
\end{proof}

\subsection{选取参数}\label{subsec:incidence-parameters}
在应用中，最大 scope $w$ 与含一个变量的因子最大数 $b$
都将被 $\log n$ 的固定幂界定。现在在该体制下选取满足
引理~\ref{lem:layout} 的参数。
计数成本的指数取决于三个量
\[
 \frac RQ\log\frac{2Q}{R},\qquad
 \tau\log(2D+1),\qquad \frac1r.
\]
它们分别对应恢复 Taylor 系数、枚举测量取值、
与在一组内枚举赋值。下面的选取同时界定三者。

\begin{lemma}\label{lem:incidence-parameters}
固定整数 $\beta\ge1$。设 $n\ge20$，
$\lambda=\lceil\log_2 n\rceil$、
$\nu=\lceil\log_2\lambda\rceil$，
并设 $b,w\le\lambda^\beta$、$n\le M\le bn$。
对所有足够大的 $n$，参数
\begin{equation}\label{eq:parameter-QR}
 r=\left\lfloor\frac{\sqrt\lambda}{\nu^2}\right\rfloor,
 \qquad Q=\left\lfloor\frac{\lambda}{128\nu}\right\rfloor,
 \qquad R=(\beta+3)\lceil\sqrt\lambda\rceil,
\end{equation}
\[
 \tau=\frac1{r\lambda},\qquad
 \epsilon=\frac1{br^2\lambda},\qquad D=r^{2Q},
\]
连同证明中指定的区间，满足引理~\ref{lem:layout}
的全部数值假设。它们还满足
\[
 R\ge32,\qquad 6R\le Q,\qquad D\le2^{\lambda/128},
\]
以及
\begin{equation}\label{eq:balanced-cost}
 \frac RQ\log\frac{2Q}{R}+\tau\log(2D+1)+\frac1r
 =O_\beta\!\left(\frac{\nu^2}{\sqrt\lambda}\right).
\end{equation}
给定数值 $n,M,b,w$，这些参数与数值假设可以在
$\poly_\beta(\log n)$ 次位运算内精确计算并检验。
对固定的 $\beta$，只有有限多个 $n$ 通不过这些检验。
\end{lemma}
\begin{proof}
对足够大的 $n$，参数的阶为
\[
 r=\Theta(\sqrt\lambda/\nu^2),\qquad
 Q=\Theta(\lambda/\nu),\qquad R=\Theta_\beta(\sqrt\lambda).
\]
特别地，最终有 $r\ge2$、$R\ge32$ 与 $6R\le Q$。
由 $\log_2 r\le\nu/2$，
\[
 \log_2 D=2Q\log_2 r
 \le2\frac{\lambda}{128\nu}\frac\nu2=\frac{\lambda}{128}.
\]
对足够大的 $n$，两个平衡不等式随之成立：
\[
 2(b+w)D^2\le4\lambda^\beta2^{\lambda/64}\le n\le M,
 \qquad 2Qr^2\le n.
\]
对第二个：$Q,r$ 被 $\lambda$ 的固定幂界定，
而 $n>2^{\lambda-1}$。第一个由同样的比较得出。

令
\[
 \mu=\frac{Q(r-1)}{2r},\qquad
 A_0=\left\lfloor\mu-\frac{R-1}{2}\right\rfloor,
 \qquad \mathcal I=\{A_0,\ldots,A_0+R-1\}.
\]
这里 $Q/4\le\mu\le Q/2$。由 $R\le Q/6$，区间落在
$\{0,\ldots,Q\}$ 内。其中心与 $\mu$ 相差小于一，
而 $R\ge32$ 保证它包含与 $\mu$ 距离在 $R/4$ 内的每个
整数。对 $W\sim\operatorname{Bin}(Q,(r-1)/(2r))$，
Hoeffding 不等式给出
\[
 \Pr[W\notin\mathcal I]
 \le2e^{-R^2/(8Q)}
 \le2e^{-16(\beta+3)^2\nu}
 \le2^{-(\beta+2)\nu}
 \le\lambda^{-(\beta+2)}
 \le\frac1{br^2\lambda}.
\]
最后一个不等式用 $b\le\lambda^\beta$ 与 $r^2\le\lambda$。
这证明了取 $\epsilon=\tau/(rb)$ 的
\eqref{eq:incidence-tail}。且一旦 $r\ge2$ 就有
$0<\tau\le1/2$。

最后，
\[
 \frac RQ=O_\beta(\nu/\sqrt\lambda),\qquad
 \log(2Q/R)=O(\nu),\qquad
 \tau\log(2D+1)=O(1/r).
\]
连同 $1/r=O(\nu^2/\sqrt\lambda)$，这些证明
\eqref{eq:balanced-cost}。

最后我们证明 \eqref{eq:parameter-QR}
中的参数可以在所述位界内精确构造。
算法找出 $\lambda=\lceil\log_2 n\rceil$
（即使 $n\le2^\lambda$ 的最小整数），
再以同样方式找出 $\nu=\lceil\log_2\lambda\rceil$。
它把 $r$ 取为满足 $r^2\nu^4\le\lambda$ 的最大整数，
用整数除法计算 $Q$，用整数平方根搜索找
$\lceil\sqrt\lambda\rceil$ 得 $R$。
若 $r<2$ 或 $Q\le R$，算法在此阶段报告失败。
否则它用整数或有理算术构造 $\tau,\epsilon,D$
与 $\mathcal I$ 的端点。然后检验
$\mathcal I\subseteq\{0,\ldots,Q\}$、两个平衡不等式、
$R\ge32$ 与 $6R\le Q$。
\eqref{eq:incidence-tail} 中的尾概率分母为 $(2r)^Q$，
故其所需界也可用整数算术检验。
这些计算花费 $\poly_\beta(\log n)$ 次位运算。
前述估计表明所有检验对足够大的 $n$ 都通过。
\end{proof}
